% Part: incompleteness
% Chapter: introduction
% Section: historical-background

\documentclass[../../../include/open-logic-section]{subfiles}

\begin{document}

\olfileid[pt]{inc}{int}{bgr}

\olsection{Contexto histórico}

Nesta seção, discutiremos brevemente alguns desenvolvimentos históricos
que ajudarão a situar os teoremas da incompletude. Em particular,
apresentaremos um panorama bastante esquemático da história da lógica
matemática e, em seguida, diremos algumas palavras sobre a história dos
fundamentos da matemática.

\begin{digress}
  A expressão ``lógica matemática'' é ambígua. Pode-se interpretar a
  palavra ``matemática'' como uma descrição do objeto de estudo, como em
  ``a lógica da matemática'', designando os princípios do raciocínio
  matemático; ou como uma descrição dos métodos, como em ``a matemática
  da lógica'', designando um estudo matemático dos princípios do
  raciocínio. O relato a seguir envolve a lógica matemática nos dois
  sentidos, muitas vezes simultaneamente.
\end{digress}

O estudo da lógica começou, essencialmente, com Aristóteles, que viveu
aproximadamente entre 384 e 322 \textsc{a.e.c.}. Suas obras
\emph{Categorias}, \emph{Analíticos anteriores} e \emph{Analíticos
  posteriores} incluem estudos sistemáticos dos princípios do
raciocínio científico, entre eles um estudo minucioso e sistemático do
silogismo.

A lógica de Aristóteles dominou a filosofia escolástica durante toda a
Idade Média; de fato, ainda no século XVIII, Kant sustentava que a
lógica aristotélica era perfeita e não precisava de revisão. Mas a
teoria do silogismo é limitada demais para modelar mais do que os
aspectos mais superficiais do raciocínio matemático. Um século antes,
Leibniz, contemporâneo de Newton, imaginara um ``cálculo'' completo para
o raciocínio lógico e dera alguns passos rudimentares na concepção de
tal cálculo, descrevendo essencialmente uma versão da lógica
proposicional.

O século XIX foi um divisor de águas para a lógica. Em 1854, George
Boole escreveu \emph{As leis do pensamento}, contendo um estudo
algébrico minucioso da lógica proposicional que não está muito distante
das apresentações modernas. Em 1879, Gottlob Frege publicou sua
\emph{Begriffsschrift} (Conceitografia), que estende a lógica
proposicional com quantificadores e relações e, assim, inclui a lógica
de primeira ordem. Na verdade, os sistemas lógicos de Frege incluíam
também a lógica de ordem superior e muito mais. Em suas \emph{Leis
  básicas da aritmética}, Frege procurou mostrar que toda a aritmética
poderia ser derivada, em sua Begriffsschrift, de pressupostos puramente
lógicos. Infelizmente, esses pressupostos revelaram-se inconsistentes,
como Russell mostrou em 1902. Mas, deixando de lado o axioma
inconsistente, Frege praticamente inventou sozinho a lógica moderna,
uma realização extraordinária. A lógica quantificacional também foi
desenvolvida de maneira independente por pensadores de orientação
algébrica posteriores a Boole, entre eles Peirce e Schr\"oder.

Voltemos agora aos desenvolvimentos nos fundamentos da matemática.
Naturalmente, como a lógica desempenha um papel importante na
matemática, há muita interação com os desenvolvimentos que acabamos de
descrever. Frege, por exemplo, desenvolveu sua lógica com o propósito
explícito de mostrar que toda a matemática poderia basear-se unicamente
em seu arcabouço lógico; em particular, desejava mostrar que a
matemática consiste em verdades \emph{analíticas} a priori, e não em
verdades \emph{sintéticas} a priori, como sustentara Kant.

Muitos consideram que o nascimento da matemática propriamente dita
ocorreu com os gregos. Os \emph{Elementos} de Euclides, escritos por
volta de 300 a.e.c., já são um representante maduro da matemática
grega, com sua ênfase no rigor e na precisão. As definições e provas
dos \emph{Elementos} de Euclides sobrevivem mais ou menos intactas nos
livros didáticos atuais de geometria do ensino médio (na medida em que
a geometria ainda é ensinada nesse nível). Esse modelo de raciocínio
matemático tem sido considerado um paradigma de argumentação rigorosa
não apenas na matemática, mas também em ramos da filosofia. (Espinosa
chegou a apresentar argumentos morais e religiosos no estilo
euclidiano, o que é curioso de ver!)

O cálculo foi inventado por Newton e Leibniz no século XVII. (Uma
acirrada disputa de prioridade prolongou-se por séculos, mas hoje a
maioria dos estudiosos considera que os dois desenvolvimentos foram,
em grande medida, independentes.) O cálculo envolve raciocínios sobre,
por exemplo, somas infinitas de quantidades infinitamente pequenas;
essas características alimentaram as críticas do bispo Berkeley, que
argumentava que a crença em Deus não era menos racional do que a
matemática de sua época. Os métodos do cálculo foram amplamente usados
no século XVIII, por exemplo por Leonhard Euler, que empregou cálculos
com somas infinitas para obter resultados impressionantes.

No século XIX, os matemáticos procuraram responder às críticas de
Berkeley assentando o cálculo sobre fundamentos mais firmes. Os
esforços de Cauchy, Weierstrass, Bolzano e outros levaram às nossas
definições atuais de limite, continuidade, diferenciação e integração
em termos de ``épsilons e deltas'', isto é, sem qualquer referência a
infinitésimos. Mais tarde naquele século, os matemáticos procuraram ir
além e explicar todos os aspectos do cálculo, inclusive os próprios
números reais, em termos dos números naturais. (Kronecker: ``Deus criou
os números inteiros; todo o resto é obra do homem.'') Em 1872, Dedekind
escreveu ``Continuidade e números irracionais'', no qual mostrou como
``construir'' os números reais como conjuntos de números racionais
(que, como você sabe, podem ser vistos como pares de números naturais);
em 1888, escreveu ``Was sind und was sollen die Zahlen'' (aproximadamente,
``O que são e para que servem os números?''), que buscava explicar os
números naturais em termos puramente ``lógicos''. Em 1887, Kronecker
escreveu ``\"Uber den Zahlbegriff'' (``Sobre o conceito de número''),
no qual falou em representar todo objeto matemático em termos dos
inteiros; em 1889, Giuseppe Peano apresentou axiomas formais e
simbólicos para os números naturais.

O fim do século XIX trouxe também uma nova ousadia no tratamento do
infinito. Antes disso, objetos e estruturas infinitários (como o
conjunto dos números naturais) eram tratados com cautela; ``infinitos''
era entendido como ``tantos quantos se queira'', e ``aproxima-se no
limite'' como ``fica tão próximo quanto se queira''. Mas Georg Cantor
mostrou que era possível tomar o infinito em sentido literal. Os
trabalhos de Cantor, Dedekind e outros ajudaram a introduzir a
compreensão geral da matemática em termos da teoria dos conjuntos, hoje
amplamente aceita.

Isso nos leva aos desenvolvimentos do século XX na lógica e nos
fundamentos. Em 1902, Russell descobriu o paradoxo no sistema lógico de
Frege. Em 1904, Zermelo provou o princípio da boa ordenação de Cantor
usando o chamado ``axioma da escolha''; a legitimidade desse axioma
provocou intenso debate. Entre 1910 e 1913, apareceram os três volumes
dos \emph{Principia Mathematica} de Russell e Whitehead, ampliando o
programa fregeano de estabelecer a matemática sobre fundamentos
lógicos. Infelizmente, Russell e Whitehead foram obrigados a adotar
dois princípios que pareciam difíceis de justificar como puramente
lógicos: um axioma do infinito e um axioma da ``redutibilidade''. Na
primeira década do século XX, Poincar\'e criticou o uso de ``definições
impredicativas'' na matemática e, na década de 1910, Brouwer começou a
propor a refundação de toda a matemática sobre uma base
``intuicionista'', que evitava o uso da lei do terceiro excluído
($!A \lor \lnot !A$).

Tempos realmente estranhos!{} O programa de reduzir toda a matemática
à lógica é hoje chamado de ``logicismo'' e, em geral, considera-se que
fracassou por causa das dificuldades mencionadas acima. O programa de
desenvolver a matemática em termos de construções mentais
intuicionistas chama-se ``intuicionismo'' e é visto como impondo
restrições excessivamente severas à matemática cotidiana. Por volta da
virada do século, David Hilbert, um dos matemáticos mais influentes de
todos os tempos, era um firme defensor dos novos métodos abstratos
introduzidos por Cantor e Dedekind: ``ninguém nos expulsará do paraíso
que Cantor criou para nós''. Ao mesmo tempo, era sensível às críticas
fundacionais a esses novos métodos (que, curiosamente, hoje são
chamados de ``clássicos''). Ele propôs uma maneira de conciliar as duas
coisas:
\begin{enumerate}
\item Representar os métodos clássicos por axiomas e regras formais;
  representar as questões matemáticas como !!{formula}s em um sistema
  axiomático.
\item Usar métodos seguros, ``finitários'', para provar que esses
  sistemas dedutivos formais são consistentes.
\end{enumerate}

O trabalho de Hilbert avançou muito na realização do primeiro objetivo.
Em 1899, ele já o fizera para a geometria em seu célebre livro
\emph{Fundamentos da geometria}. Nos anos seguintes, ele e vários de
seus alunos e colaboradores trabalharam em outras áreas da matemática
para fazer o que Hilbert fizera pela geometria. O próprio Hilbert
apresentou sistemas de axiomas para a aritmética e a análise. Zermelo
apresentou uma axiomatização da teoria dos conjuntos, posteriormente
ampliada por Fraenkel, Skolem, von Neumann e outros. Em meados da década
de 1920, duas abordagens reivindicavam o título de axiomatização de
``toda'' a matemática: os \emph{Principia mathematica} de Russell e
Whitehead e o que veio a ser conhecido como teoria dos conjuntos de
Zermelo--Fraenkel.

Em 1921, Hilbert iniciou um projeto de pesquisa com o objetivo de provar
que esses sistemas eram consistentes. Nesse projeto, contou com a ajuda
de vários de seus alunos, em particular Bernays, Ackermann e, mais
tarde, Gentzen. A ideia básica para alcançar esse objetivo era
formular a questão da possibilidade de uma !!{derivation} de uma
inconsistência na matemática como um problema combinatório sobre
possíveis sequências de símbolos, a saber, possíveis sequências de
sentenças que satisfizessem o critério de constituir uma
!!{derivation} correta, digamos, de $!A \land \lnot !A$ a partir dos
axiomas de um sistema axiomático para a aritmética, a análise ou a
teoria dos conjuntos. Uma prova da impossibilidade de tal sequência de
símbolos seria---por ser ela própria uma prova matemática---formalizável
nesses sistemas axiomáticos. Em outras palavras, haveria alguma
sentença $\OCon$ que afirmasse, por exemplo, que a aritmética é
consistente. Além disso, essa sentença deveria ser demonstrável nos
sistemas em questão, sobretudo se sua prova exigisse apenas meios muito
restritos, ``finitários''.

O segundo objetivo, o de que os sistemas de axiomas desenvolvidos
decidissem toda questão matemática, pode ser formulado com precisão de
duas maneiras. A primeira é a seguinte: para qualquer sentença~$!A$ na
linguagem de um sistema axiomático para a matemática, ou $!A$ ou
$\lnot !A$ é demonstrável a partir dos axiomas. Se isso fosse verdade,
não haveria sentenças que não pudessem ser nem provadas nem refutadas
com base nos axiomas, nem questões que os axiomas não decidissem. Um
sistema axiomático com essa propriedade é chamado de \emph{completo}.
Naturalmente, para uma sentença dada, ainda poderia ser difícil
determinar qual das duas alternativas se verifica. Mas, em princípio,
deveria haver um método para fazê-lo. De fato, para os sistemas de
axiomas e de !!{derivation} considerados por Hilbert, a completude
implicaria a existência de tal método---embora Hilbert não tivesse
percebido isso. A segunda maneira de interpretar a questão seria esta
exigência mais forte: que houvesse um método mecânico, computacional,
capaz de determinar, para uma dada sentença~$!A$, se ela é ou não
derivável dos axiomas.

Em 1931, G\"odel provou os dois ``teoremas da incompletude'', que
mostraram que esse programa não poderia ter êxito. Não há sistema de
axiomas completo para a matemática; em particular, a sentença que
expressa a consistência dos axiomas é uma sentença que não pode ser nem
provada nem refutada.

Isso desferiu um golpe fatal no programa original de Hilbert. Contudo,
como tantas vezes acontece na matemática, também abriu novos e
estimulantes caminhos de pesquisa. Se não existe um único sistema
formal de matemática que abarque tudo, faz sentido desenvolver sistemas
mais circunscritos e investigar o que pode ser provado neles. Também
faz sentido desenvolver métodos de prova menos restritos para
estabelecer a consistência desses sistemas e encontrar maneiras de
medir a dificuldade de provar sua consistência. Como G\"odel mostrou
que (quase) todo sistema formal tem questões que não consegue decidir,
faz sentido procurar questões ``interessantes'' que um dado sistema
formal não possa decidir e determinar quão forte deve ser um sistema
formal para decidi-las. Até hoje, os lógicos vêm investigando essas
questões em uma nova disciplina matemática, a teoria da prova.

\end{document}
